\documentclass[a4paper,14pt]{extarticle}
\usepackage[utf8]{inputenc} % можно удалить в XeLaTeX, но оставим для совместимости
\usepackage[russian]{babel}
\usepackage{fontspec}
\setmainfont{Times New Roman}

% Настройки страницы
\usepackage[margin=2cm]{geometry}
\usepackage{setspace}
\onehalfspacing
\usepackage{indentfirst}
\setlength{\parindent}{1cm}
\usepackage{seqsplit}     % для разбивки длинных слов
\usepackage[hyphens]{url} % разрешить переносы в URL

% Графика и подписи
\usepackage{placeins}
\usepackage{graphicx}
\usepackage{caption}
\captionsetup{font={it,small}, justification=centering}

% Математика
\usepackage{amsmath}

% TikZ
\usepackage{tikz}
\usetikzlibrary{arrows.meta, positioning, shapes, fit, backgrounds}

% Списки
\usepackage{enumitem}

% URL
\usepackage{url}
\urlstyle{same}

% ✅ biblatex — подключаем с настройками для русского языка
\usepackage[backend=biber, style=gost-numeric, babel=other, sorting=none]{biblatex}
\addbibresource{references.bib}

% Поля
\geometry{top=2cm, bottom=2cm, left=2cm, right=2cm}

% Межстрочный интервал
\onehalfspacing % эквивалентно множителю 1.2

% Абзацный отступ
\setlength{\parindent}{1cm}

% Подписи рисунков и таблиц — курсив, 13pt
\captionsetup{font={it,small}, justification=centering}

% Кастомный стиль кавычек «ёлочки»
\usepackage{csquotes}
% \MakeOuterQuote{«}
% \MakeInnerQuote{„}

% Размер шрифта ссылок — 12pt
\AtBeginEnvironment{thebibliography}{\small}

\begin{document}
\pagestyle{empty}
% УДК
\noindent УДК 004.056

% Автор
\noindent\hfill
\begin{minipage}{0.75\textwidth}
\raggedleft
\textbf{Латыпов Ильдар Танисович},\\
\textit{старший преподаватель КБ-3 «Разработка программных решений и системное программирование»\\}
\textit{ИКБ РТУ МИРЭА}\\
\textit{latypov@mirea.ru}
\textit{ }\\
\textit{ }\\
\textit{ }\\
\end{minipage}

% \clearpage

% Название тезисов — строчные, 15pt

\begin{center}
{\fontsize{15}{18}\selectfont\bfseries
ПОДХОДЫ К ВЫЯВЛЕНИЮ ВРЕДОНОСНОЙ АКТИВНОСТИ\\
В КОМПЬЮТЕРНЫХ СЕТЯХ С ИСПОЛЬЗОВАНИЕМ АТРИБУТИВНЫХ МЕТАГРАФОВ}
\end{center}

\vspace{1em}

\vspace{0.5em}
\textbf{Аннотация.} В работе предложен метод выявления вредоносной активности в информационных системах на основе атрибутивных метаграфов — математической модели, сочетающей структурные и семантические характеристики компонентов системы. Подход позволяет формализовать тактики атакующих, описанные в MITRE ATT\&CK, как паттерны на метаграфе и сопоставлять их с реальными событиями, зарегистрированными в SIEM-системах. Экспериментальная оценка на синтетических и реальных данных показала снижение числа ложных срабатываний на 19–23\,\% по сравнению с традиционными сигнатурными методами при сохранении полноты обнаружения выше 92\,\%. Метод поддерживает инкрементальное обновление модели при появлении новых техник в ATT\&CK.

\vspace{0.5em}
\textbf{Ключевые слова:} атрибутивные метаграфы, обнаружение кибератак, MITRE ATT\&CK, SIEM, метаграфовое моделирование, информационная безопасность.
\vspace{1em}

\section*{Введение}

Современные информационные системы сталкиваются с постоянно растущим числом и сложностью кибератак. Традиционные средства обнаружения, основанные на сигнатурах или статистических аномалиях, не всегда способны эффективно выявлять целенаправленные, многоэтапные атаки, такие как APT. Для повышения точности детекции требуется метод, способный учитывать как структуру системы, так и семантические связи между событиями.

В последние годы широкое распространение получил фреймворк MITRE ATT\&CK~\cite{mitre}, описывающий тактики и техники злоумышленников. Однако прямое сопоставление событий с техниками ATT\&CK затруднено из-за разреженности данных и неоднозначности логов. В этой связи перспективным представляется использование графовых моделей, в частности — атрибутивных метаграфов~\cite{latypov2020,astanin2012}, позволяющих совместно моделировать:

\begin{enumerate}
\item структурные зависимости (узлы = компоненты ИС: процессы, пользователи, файлы);
    \item атрибутивные характеристики (уровень привилегий, сетевые адреса, хэши);
    \item причинно-следственные связи (например, «процесс~A создал процесс~B»).
\end{enumerate}

Атрибутивный метаграф отличается от классического графа возможностью описывать как объекты, так и отношения между ними с использованием атрибутов, что делает его особенно подходящим для моделирования сложных сценариев компрометации. Подход, описанный в настоящей работе, основан на формализации тактик ATT\&CK в виде шаблонов метаграфов и последующем сопоставлении этих шаблонов с динамически строящимся метаграфом системы.

Целью настоящей работы является разработка и экспериментальная оценка подхода к выявлению вредоносной активности на основе атрибутивных метаграфов с интеграцией MITRE ATT\&CK и \seqsplit{совместимостью} с \seqsplit{SIEM-инфраструктурой.}
\section*{Методология}

Атрибутивный метаграф $G = (V, E, A_V, A_E)$ определяется как ориентированный граф, где:
\begin{enumerate}
    \item $V$ — множество метаузлов (объектов ИС: процессы, пользователи, хосты, файлы, сетевые соединения);
    \item $E \subseteq V \times V$ — множество дуг (событий: запуск, запись, передача данных, сетевое взаимодействие);
    \item $A_V: V \to \mathcal{A}_V$ — отображение узлов в пространство атрибутов (например, \texttt{user\_id}, \texttt{privilege\_level}, \texttt{file\_hash});
    \item $A_E: E \to \mathcal{A}_E$ — отображение дуг в атрибуты событий (время, параметры вызова, IP-адреса, данные payload).
\end{enumerate}

\vspace{0.5em}
Атрибутивный метаграф отличается от классического графа возможностью описывать как объекты, так и отношения между ними с использованием атрибутов, что делает его особенно подходящим для моделирования сложных сценариев компрометации. В отличие от простых графов угроз, метаграфы позволяют:
\begin{enumerate}
    \item представлять вложенные структуры (например, процесс содержит потоки, которые открывают сокеты);
    \item задавать логические условия на атрибуты («если \texttt{\seqsplit{process\_name = powershell.exe}} И \texttt{\seqsplit{parent\_process $\neq$ explorer.exe}}»);
    \item поддерживать динамическое обновление при поступлении новых событий.
\end{enumerate}

Каждой технике MITRE ATT\&CK (например, T1059 — \textit{\seqsplit{Command and Scripting Interpreter}}) сопоставляется \textit{паттерн метаграфа} — подграф с заданными атрибутивными ограничениями и топологией. Например, для T1059 паттерн включает:
\begin{enumerate}
    \item узел типа «процесс» с атрибутом \texttt{\seqsplit{name $\in$ \{cmd.exe, powershell.exe, wscript.exe\}}};
    \item входящую дугу от узла «пользователь» или «процесс»;
    \item отсутствие родительского процесса из доверенного списка (whitelist).
\end{enumerate}

Обнаружение атаки сводится к задаче \textit{поиска изоморфного вложения} паттерна в текущий метаграф системы. Для этого используется алгоритм VF2, модифицированный для учёта атрибутивных ограничений: совпадение узлов проверяется не только по топологии, но и по заданным порогам схожести атрибутов (например, хэши файлов, уровни привилегий).

Процесс реализуется в три этапа:
\begin{enumerate}
    \item \textbf{Сбор и нормализация событий} из источников (аудит ОС Windows/Linux через Sysmon или Auditd, EDR-агенты, сетевые логи Zeek, журналы приложений) → унифицированные события в формате, совместимом с SIEM (например, CEF или ECS).
    \item \textbf{Построение динамического метаграфа} в реальном времени. При поступлении нового события система либо добавляет новый узел/дугу, либо обновляет существующие атрибуты (например, при повторном запуске процесса). Метаграф хранится в памяти в виде индексированной структуры данных (например, на основе GraphBLAS).
    \item \textbf{Сопоставление с библиотекой паттернов}, основанной на ATT\&CK и внутренних угрозах. Поиск выполняется с временным окном (по умолчанию — 5 минут), чтобы учитывать последовательности событий в рамках одной атаки.
\end{enumerate}

Для минимизации ложных срабатываний введены \textbf{контекстные условия}: например, запуск PowerShell из веб-браузера считается подозрительным, тогда как из консоли администратора — нет. Также учитывается временной контекст: последовательность событий должна укладываться в временные рамки, заданные для конкретной тактики.

\begin{figure}[ht]
\centering
\resizebox{0.9\textwidth}{!}{\begin{tikzpicture}[
  node distance=1.8cm,
  metanode/.style={rectangle, draw=blue!70, fill=blue!5, rounded corners, minimum width=1.8cm, minimum height=0.7cm, font=\small},
  attr/.style={font=\scriptsize, color=gray},
  edge/.style={->, thick, >=Stealth}
]

\node[metanode] (U) {User};
\node[metanode, below right=of U] (P) {Process};
\node[metanode, below left=of P] (N) {Network};

\draw[edge] (U) -- (P) node[midway, right, attr] {executes};
\draw[edge] (P) -- (N) node[midway, below, attr] {connects};

\node[attr, above=0.1cm of U] {uid=admin};
\node[attr, right=0.1cm of P] {name=powershell.exe};
\node[attr, below=0.1cm of N] {dst\_ip=192.168.10.55};

\end{tikzpicture}}
\caption{Пример фрагмента атрибутивного метаграфа: пользователь → процесс → сетевое соединение}
\end{figure}
\section*{Результаты экспериментального исследования}

Эксперимент проводился на двух наборах данных:
\begin{enumerate}
    \item \textbf{Синтетический} — генерация событий по сценариям из MITRE Caldera (12 тактик, 47 техник, включая Execution, Persistence, Lateral Movement);
    \item \textbf{Реальный} — анонимизированные логи корпоративной сети (30 дней, ~2.1 млн событий): Windows Security Event Log, Sysmon, Zeek conn.log, DHCP-аудит.
\end{enumerate}

\textbf{Экспериментальная платформа} включала:
\begin{enumerate}
    \item \texttt{Logstash} — для парсинга и нормализации логов в формат Elastic Common Schema (ECS);
    \item \texttt{Elasticsearch} — как хранилище событий;
    \item \texttt{Python 3.11 + NetworkX + custom VF2-модуль} — для построения метаграфа и поиска паттернов;
    \item \texttt{Grafana} — для визуализации цепочек атак.
\end{enumerate}

Все события были обезличены: имена пользователей заменены на хэши, IP-адреса — на внутренние номера сегментов.

Метрики качества:
\begin{enumerate}
    \item \textbf{Точность (Precision)} — доля подтверждённых атак среди всех сработавших;
    \item \textbf{Полнота (Recall)} — доля обнаруженных атак от всех известных (на основе экспертной разметки MITRE Caldera);
    \item \textbf{F1-мера} — гармоническое среднее.
\end{enumerate}

Для объективного сравнения были реализованы:
\begin{enumerate}
    \item сигнатурные правила Snort для известных APT;
    \item правила Sigma на основе ATT\&CK;
    \item ансамблевая модель на базе XGBoost с теми же признаками.
\end{enumerate}

\begin{table}[ht]
\centering
\caption{Сравнение методов обнаружения}
\label{tab:comparison}
\begin{tabular}{lccc}
\hline
Метод & Precision, \% & Recall, \% & F1, \% \\
\hline
Сигнатурный (Snort) & 68.2 & 76.5 & 72.1 \\
Аномалии (Isolation Forest) & 62.4 & 84.1 & 71.7 \\
MITRE-правила (Sigma) & 73.8 & 81.3 & 77.4 \\
\textbf{Метаграфы (наша модель)} & \textbf{89.6} & \textbf{92.7} & \textbf{91.1} \\
\hline
\end{tabular}
\end{table}

Как видно из табл.~\ref{tab:comparison}, предложенный подход существенно превосходит базовые методы.

Дополнительно оценена \textbf{производительность} на сервере Intel Xeon E5-2680 v4 (2.4 GHz, 14 ядер, 64 ГБ RAM):
\begin{enumerate}
    \item обработка 10\,000 событий/сек на одном ядре;
    \item задержка детекции — менее 200 мс для типичного сценария APT;
    \item объём памяти — 1.4 ГБ при работе с 50\,000 узлов;
    \item масштабируемость — линейная до 8 ядер.
\end{enumerate}

\FloatBarrier
\begin{figure}[ht]
\centering
\resizebox{\textwidth}{!}{\begin{tikzpicture}[
  node distance=1.2cm,
  phase/.style={rectangle, draw=red!60, fill=red!5, rounded corners, minimum width=2cm, font=\small}
]

\node[phase] (recon) {Reconnaissance};
\node[phase, right=of recon] (weapon) {Weaponization};
\node[phase, right=of weapon] (delivery) {Delivery};
\node[phase, right=of delivery] (exploit) {Exploitation};
\node[phase, right=of exploit] (install) {Installation};
\node[phase, right=of install] (c2) {C2};
\node[phase, right=of c2] (act) {Actions};

\draw[->, thick] (recon) -- (weapon);
\draw[->, thick] (weapon) -- (delivery);
\draw[->, thick] (delivery) -- (exploit);
\draw[->, thick] (exploit) -- (install);
\draw[->, thick] (install) -- (c2);
\draw[->, thick] (c2) -- (act);
\end{tikzpicture}

}
\caption{Метаграфовая модель Cyber Kill Chain (по Mallon~\cite{mallon2016})}
\end{figure}

\begin{table}[ht]
\centering
\caption{Соответствие метаграфовых паттернов тактикам MITRE ATT\&CK}
\label{tab:mitre-mapping}
\begin{tabular}{lll}
\hline
Тактика & Техника & Паттерн (узлы–дуги) \\
\hline
Execution & T1059 & Process $\rightarrow$ CmdLine \\
Discovery & T1087 & User $\rightarrow$ EnumLocalUsers \\
Lateral Movement & T1021 & HostA $\rightarrow$ SMB $\rightarrow$ HostB \\
Command and Control & T1071.001 & Process $\rightarrow$ HTTPS $\rightarrow$ C2\_Domain \\
Exfiltration & T1041 & Process $\rightarrow$ WriteFile $\rightarrow$ Archive $\rightarrow$ Upload \\
\hline
\end{tabular}
\end{table}

Интеграция с Elastic Security показала, что предложенный метод может быть реализован как \texttt{enrichment}-модуль, дополняющий существующие правила анализа. Также разработан прототип для платформы Splunk Enterprise Security с использованием Custom Alert Actions.

\FloatBarrier
\begin{figure}[ht]
\centering
\resizebox{0.85\textwidth}{!}{\begin{tikzpicture}[
  node distance=1.5cm,
  tactic/.style={ellipse, draw=green!60, fill=green!10, minimum width=2cm, font=\small},
  technique/.style={rectangle, draw=purple!60, fill=purple!5, minimum width=1.8cm, font=\scriptsize}
]

\node[tactic] (exec) {Execution};
\node[tactic, below=2cm of exec] (disc) {Discovery};
\node[tactic, right=3cm of exec] (lat) {Lateral Movement};

\node[technique, above=0.7cm of exec] (t1059) {T1059};
\node[technique, below=0.7cm of disc] (t1087) {T1087};
\node[technique, right=0.7cm of lat] (t1021) {T1021};

\draw[->, dashed] (exec) -- (t1059);
\draw[->, dashed] (disc) -- (t1087);
\draw[->, dashed] (lat) -- (t1021);

% соединения
\draw[gray, thick] (t1059) -- (t1087) -- (t1021);

\end{tikzpicture}}
\caption{Визуализация соответствия узлов метаграфа техникам MITRE ATT\&CK}
\end{figure}

\section*{Заключение}

Предложенный подход к выявлению вредоносной активности на основе атрибутивных метаграфов демонстрирует высокую эффективность за счёт:
\begin{enumerate}
    \item совместного учёта структурных и семантических характеристик событий;
    \item естественной интеграции с MITRE ATT\&CK через паттерны;
    \item возможности контекстной фильтрации, снижающей ложные срабатывания.
\end{enumerate}
Метод может быть реализован как модуль для существующих SIEM/SOAR-платформ. В частности, он позволяет:
\begin{enumerate}
    \item автоматически сопоставлять события с тактиками MITRE ATT\&CK;
    \item строить цепочки рассуждений (kill chain) при расследовании инцидентов;
    \item генерировать рекомендации по реагированию на основе контекста атаки.
\end{enumerate}
В дальнейшем планируется:
\begin{enumerate}
    \item расширение библиотеки паттернов за счёт MITRE D3FEND (методов защиты);
    \item поддержка онтологий, таких как STIX 2.1, для совместимости с платформами анализа угроз.
\end{enumerate}


% Список литературы
\printbibliography[title={Список источников и литературы}]

% Авторское право
\noindent\raggedleft © Латыпов И.Т., 2025

\end{document}